\section{Uso en frases cotidianas}
\textbf{Frases modelo}
\begin{itemize}[leftmargin=*]
  \item \textit{Kan u me helpen?} (¿Puede ayudarme?)
  \item \textit{Kan je langzamer spreken?} (¿Puedes hablar más despacio?)
  \item \textit{Ik moet nu vertrekken.} (Tengo que irme ahora.)
  \item \textit{We moeten boodschappen doen}. (Tenemos que hacer compras.)
  \item \textit{Ik wil graag een broodje kaas}. (Quiero un bocadillo de queso, cortés.)
  \item \textit{Wil je iets drinken?} (¿Quieres beber algo?)
\end{itemize}

\textbf{Diálogo corto}
\begin{itemize}[leftmargin=*]
  \item A:\@ \textit{Kun je morgen komen?}
  \item B:\@ \textit{Nee, ik moet werken. Maar ik wil vrijdag afspreken.}
\end{itemize}

\textbf{Mini práctica}
\begin{itemize}[leftmargin=*]
  \item Escribe dos frases de permiso/capacidad con \textit{kan}.
  \item Redacta una excusa con \textit{moet} y una invitación con \textit{wil}.
\end{itemize}
